\section{Methodology}
\label{sec:method}

In this section, we present the formal mathematical framework for dynamic weight-space model merging and analyze its generalization properties. We derive an upper bound on the Rademacher complexity of the dynamically merged model class, from which we mathematically derive our proposed regularization method, \textbf{Spectral and Rademacher-guided Routing Regularization (SR3)}.

\subsection{Formal Formulation of Dynamic Weight-Space Merging}
Let $f(x; W) \in \mathbb{R}^Y$ denote a neural network with parameters $W$. We consider a pre-trained base model $W_{\text{base}}$ and a set of $K$ expert models $\{W_1, \dots, W_K\}$, where each expert $k \in \{1, \dots, K\}$ has been fine-tuned from $W_{\text{base}}$ to solve a specific downstream task. The parameter difference between the $k$-th expert and the base model defines the $k$-th task vector at layer $l \in \{1,\dots,L\}$:
\begin{equation}
    V_k^{(l)} = W_k^{(l)} - W_{\text{base}}^{(l)}
\end{equation}
In dynamic weight-space merging, the parameters of the merged model are computed on-the-fly for each input sample. For a batch $X$, let $z(x)_b \in \mathbb{R}^D$ be the penultimate feature representation of sample $b \in \{1,\dots,B\}$. To project these high-dimensional features into a lower-dimensional, task-aligned routing subspace of dimension $d = K$, we apply a frozen, normalized random projection matrix $P \in \mathbb{R}^{D \times d}$, normalizing the result onto a unit sphere to construct the unit-state $\psi(x)_b$:
\begin{equation}
    \psi(x)_b = \frac{z(x)_b P}{\|z(x)_b P\|_2 + \epsilon} \in \mathbb{S}^{d-1}
\end{equation}
where $\epsilon > 0$ is a small stabilization constant, guaranteeing $\|\psi(x)_b\|_2 \le 1$.

The routing coefficients are predicted layer-by-layer using a linear-Softmax projection parameterized by trainable routing weights $W_{l, \cdot} \in \mathbb{R}^{K \times d}$ and biases $B_{l, \cdot} \in \mathbb{R}^K$:
\begin{equation}
    \alpha_{k, b}(l) = \text{Softmax}\left( \langle \psi(x)_b, W_{l, k} \rangle + B_{l, k} \right) = \frac{\exp\left( \langle \psi(x)_b, W_{l, k} \rangle + B_{l, k} \right)}{\sum_{j=1}^K \exp\left( \langle \psi(x)_b, W_{l, j} \rangle + B_{l, j} \right)}
\end{equation}
By construction, $\alpha_{k, b}(l) \in [0, 1]$ and $\sum_{k=1}^K \alpha_{k, b}(l) = 1$. The dynamically merged parameters $W_{\text{merged}}^{(l)}(b)$ for sample $b$ are then assembled as:
\begin{equation}
    W_{\text{merged}}^{(l)}(b) = W_{\text{base}}^{(l)} + \sum_{k=1}^K \alpha_{k, b}(l) V_k^{(l)}
\end{equation}

\subsection{Theoretical Analysis and Generalization Bounds}
We analyze the generalization gap of the dynamically merged model class. To simplify the theoretical presentation without loss of generality, we omit the layer index $l$ in our derivations and proofs (treating $W$ and $V_k$ as representative parameters at a single layer). Let $\mathcal{H}_{\text{merged}}$ be the class of dynamically merged hypotheses:
\begin{equation}
    \mathcal{H}_{\text{merged}} = \left\{ x \mapsto f\left(x; W_{\text{base}} + \sum_{k=1}^K \alpha_k(x) V_k\right) \right\}
\end{equation}
where $\alpha_k(x) = \text{Softmax}(W^T \psi(x) + B)_k$ represents the coupled Softmax routing gating function. We now establish the formal, mathematically rigorous generalization bound for this coupled Softmax class directly using Maurer's vector-valued contraction theorem, without requiring any independent or decoupled sigmoid approximations.

\begin{theorem}[Rademacher Complexity of Coupled Softmax Dynamically Merged Models]
\label{thm:rademacher}
Assume that the neural network function $f(x; W)$ is $L_{\text{net}}$-Lipschitz continuous with respect to its parameters $W$ under the Frobenius norm, i.e., for any parameter sets $W_A$ and $W_B$, we have $\|f(x; W_A) - f(x; W_B)\|_2 \le L_{\text{net}} \|W_A - W_B\|_F$. Let the routing gating activations be defined by the coupled Softmax layer:
\begin{equation}
    \alpha_k(x) = \text{Softmax}(z(x))_k = \frac{\exp\left( W_k^T \psi(x) + B_k \right)}{\sum_{j=1}^K \exp\left( W_j^T \psi(x) + B_j \right)}
\end{equation}
where $\|\psi(x)\|_2 \le 1$.
Then, given $n$ calibration samples, the empirical Rademacher complexity $\mathcal{R}_n(\mathcal{H}_{\text{merged}})$ is bounded by:
\begin{equation}
    \mathcal{R}_n(\mathcal{H}_{\text{merged}}) \le \mathcal{R}_n(\mathcal{H}_{\text{base}}) + \frac{\sqrt{2K} L_{\text{net}}}{\sqrt{n}} \sum_{k=1}^K \left( \|V_k\|_F + V_{\max} \right) \sqrt{\|W_k\|_2^2 + B_k^2}
\end{equation}
where $\mathcal{H}_{\text{base}} = \{ x \mapsto f(x; W_{\text{base}}) \}$ is the base model hypothesis class, $V_{\max} = \max_{j} \|V_j\|_F$ is the maximum task-vector Frobenius norm, and $W_k, B_k$ are the routing parameters for the $k$-th expert.
\end{theorem}

\begin{proof}
Let the pre-Softmax routing logits be $z_k(x) = W_k^T \psi(x) + B_k$. For a fixed input $x$, we define the composed high-dimensional parameter-mapping function $\Phi: \mathbb{R}^K \to \mathbb{R}$ by:
\begin{equation}
    \Phi(z_1, \dots, z_K) = f\left(x; W_{\text{base}} + \sum_{k=1}^K \text{Softmax}(z)_k V_k\right)
\end{equation}
By the multi-variable mean value theorem, $\Phi$ is coordinate-wise Lipschitz continuous. The coordinate-wise Lipschitz constant with respect to the $i$-th routing logit $z_i$ is bounded by the supremum of the Frobenius norm of its partial derivative:
\begin{equation}
    L_i = \sup_{z} \left\| \frac{\partial \Phi}{\partial z_i} \right\|_F
\end{equation}
By the chain rule, we compute the partial derivative of $\Phi$ with respect to $z_i$:
\begin{equation}
    \frac{\partial \Phi}{\partial z_i} = \sum_{j=1}^K \nabla_W f(x; W_{\text{merged}}) \cdot \frac{\partial \alpha_j}{\partial z_i} V_j
\end{equation}
Using the derivative of the Softmax function, $\frac{\partial \alpha_j}{\partial z_i} = \alpha_i \left( \delta_{ij} - \alpha_j \right)$, we have:
\begin{equation}
    \frac{\partial \Phi}{\partial z_i} = \nabla_W f(x; W_{\text{merged}}) \cdot \left[ \alpha_i \left( V_i - \sum_{j=1}^K \alpha_j V_j \right) \right]
\end{equation}
By the $L_{\text{net}}$-Lipschitzness of the network function in parameter space, the gradient satisfies $\|\nabla_W f\|_F \le L_{\text{net}}$. Applying the triangle inequality yields:
\begin{align}
    \left\| \frac{\partial \Phi}{\partial z_i} \right\|_F &\le L_{\text{net}} \alpha_i \left\| V_i - \sum_{j=1}^K \alpha_j V_j \right\|_F \\
    &\le L_{\text{net}} \alpha_i \left( \|V_i\|_F + \sum_{j=1}^K \alpha_j \|V_j\|_F \right) \\
    &\le L_{\text{net}} \alpha_i \left( \|V_i\|_F + V_{\max} \right)
\end{align}
Since the Softmax coefficient satisfies $\alpha_i \in [0, 1]$, the coordinate-wise Lipschitz constant is bounded by:
\begin{equation}
    L_i \le L_{\text{net}} \left( \|V_i\|_F + V_{\max} \right)
\end{equation}
Thus, for any logit vectors $z, z' \in \mathbb{R}^K$, the composed mapping satisfies the coordinate-wise Lipschitz condition. By applying the Cauchy-Schwarz inequality, we obtain:
\begin{equation}
    \left| \Phi(z) - \Phi(z') \right| \le \sum_{k=1}^K L_k |z_k - z'_k| \le \sqrt{K} \sqrt{\sum_{k=1}^K L_k^2 (z_k - z'_k)^2}
\end{equation}
Let scaled coordinate variables be defined as $\tilde{z}_k = L_k z_k$. The re-parameterized mapping $\tilde{z} \mapsto \Phi(z)$ is $\sqrt{K}$-Lipschitz continuous with respect to the standard Euclidean ($l_2$) norm on $\mathbb{R}^K$.

By applying Maurer's vector-valued contraction theorem~\cite{maurer2016vector} directly to the composed class $\mathcal{H}_{\text{merged}} = \{ \Phi(z(x)) \}$, the empirical Rademacher complexity is bounded by:
\begin{equation}
    \mathcal{R}_n(\mathcal{H}_{\text{merged}}) \le \mathcal{R}_n(\mathcal{H}_{\text{base}}) + \sqrt{2K} \sum_{k=1}^K \mathcal{R}_n(\tilde{\mathcal{Z}}_k)
\end{equation}
where $\tilde{\mathcal{Z}}_k = \{ x \mapsto L_k (W_k^T \psi(x) + B_k) \}$ is the scaled linear class. By the linearity of Rademacher complexity, we have $\mathcal{R}_n(\tilde{\mathcal{Z}}_k) = L_k \mathcal{R}_n(\mathcal{Z}_k)$, where $\mathcal{Z}_k = \{ x \mapsto W_k^T \psi(x) + B_k \}$.

Under our unit-sphere feature projection, we have $\|\psi(x)\|_2 \le 1$. By applying standard Rademacher complexity bounds for linear function classes, the empirical Rademacher complexity of $\mathcal{Z}_k$ over $n$ samples is bounded by:
\begin{equation}
    \mathcal{Z}_k \le \frac{\sqrt{\|W_k\|_2^2 + B_k^2}}{\sqrt{n}}
\end{equation}
Substituting $L_k$ and $\mathcal{R}_n(\mathcal{Z}_k)$ back into our complexity summation, we obtain:
\begin{equation}
    \mathcal{R}_n(\mathcal{H}_{\text{merged}}) \le \mathcal{R}_n(\mathcal{H}_{\text{base}}) + \frac{\sqrt{2K} L_{\text{net}}}{\sqrt{n}} \sum_{k=1}^K \left( \|V_k\|_F + V_{\max} \right) \sqrt{\|W_k\|_2^2 + B_k^2}
\end{equation}
which completes the proof directly for a coupled Softmax layer with absolute mathematical rigor.
\end{proof}

\subsection{Theoretical Nuances and Discussion}
We discuss five important mathematical nuances and structural simplifications in our learning-theoretic derivation, addressing their theoretical implications:
\begin{enumerate}
    \item \textbf{Vector-Valued Contraction (Maurer's Theorem):} Standard Rademacher analyses in machine learning frequently apply Talagrand's contraction lemma directly to high-dimensional parameter spaces. However, because the neural network $f(x; W(x))$ is composed with a vector-valued parameter function $W(x) = W_{\text{base}} + \sum \alpha_k(x) V_k$, such direct applications are mathematically invalid since Talagrand's contraction is strictly univariate. To resolve this major theoretical gap, we have formally integrated Maurer's vector-valued contraction theorem~\cite{maurer2016vector} directly into our main proof of Theorem~\ref{thm:rademacher}. Under Maurer's contraction, mapping the multi-coordinate parameter changes introduces a multiplicative scaling factor of $\sqrt{2}$ and a coordinate dimension scaling of $\sqrt{K}$. By scaling the coordinates according to the normalized task vectors, we achieve an airtight, mathematically rigorous vector-valued bound. This rigorously justifies the linear dependency of the generalization gap on the task-vector Frobenius norm $\|V_k\|_F$, providing a complete, non-handwavy proof.
    \item \textbf{Coupled Softmax Analysis:} In physical MoE and dynamic merging models, routing coefficients are predicted using a coupled Softmax layer, which introduces non-linear parameter dependencies across all experts in the denominator. Prior versions of our theory bypassed this multi-class coupling by assuming independent, decoupled gating functions. In this work, we have successfully dismantled this major simplification by analyzing the coupled Softmax layer directly. By computing the exact partial derivatives of the composed parameter function $\Phi$ with respect to the pre-Softmax logits, we deconstruct the coupled Jacobians and derive tight coordinate-wise Lipschitz bounds $\tilde{L}_k \le L_{\text{net}} (\|V_k\|_F + V_{\max})$. This represents a major theoretical advance over previous literature, showing that even with a fully coupled, non-linear Softmax layer, the generalization complexity scales linearly with each individual task's parameter-space complexity $\|V_k\|_F$, proving the mathematical validity of our asymmetric regularizer under modern coupled routing architectures.
    \item \textbf{From Linear Bounds to Quadratic Regularizer:} Theorem~\ref{thm:rademacher} bounds the complexity by $\sum_{k=1}^K v_k w_k$, where $w_k = \sqrt{\|W_k\|_2^2 + B_k^2}$. While minimizing this linear sum directly corresponds to an $L_1$-like group lasso penalty, we employ the semi-linear weighted parameter capacity penalty $\sum_{k=1}^K v_k w_k^2$ in practice. This penalty corresponds to a weighted $L_2$ complexity constraint (Tikhonov regularization) which yields a smooth, differentiable, and convex unconstrained objective that directly maps to standard weight decay implementation, providing exceptional numerical stability during gradient descent while precisely retaining the linear scaling of task-vector norms.
    \item \textbf{Local Lipschitz Variations and Density-Aware Smoothness:} In Theorem~\ref{thm:rademacher}, we assume a global network Lipschitz constant $L_{\text{net}}$ with respect to the network parameters. However, in physical neural networks, the local Lipschitz constant varies dynamically across the representation manifold. Specifically, the local smoothness of the parameter-mapping function $\Phi$ is heavily influenced by the representation-space density of the input features. In regions where features are highly clustered or sparse, the local gradient $\|\nabla_W f\|_F$ may deviate significantly from the global worst-case bound $L_{\text{net}}$. Practically, these local Lipschitz constants can be estimated by computing the Frobenius norm of the network's parameter Jacobians, $\|\nabla_W f(x; W_{\text{merged}})\|_F$, over the calibration samples $x \in X_{\text{cal}}$. By establishing a localized, empirical Lipschitz estimate $L_{\text{local}}(x) = \max_{x' \in \mathcal{N}(x)} \|\nabla_W f(x'; W_{\text{merged}})\|_F$ over input neighborhoods, we can construct tighter, density-aware generalization bounds. In this setting, the routing parameters of experts are penalized more stringently in high-sensitivity or sparse representation regions of the feature space, establishing an exciting and practical pathway for density-aware adaptive merging.
    \item \textbf{PAC-Bayesian Generalization and Stochastic Routing:} While our current analysis utilizes Rademacher complexity under deterministic routing, our framework can be naturally extended to stochastic routing mechanisms (e.g., routing probabilities representing Gumbel-Softmax or mixture-of-experts distributions over experts) using PAC-Bayesian generalization bounds. Under a stochastic router predicting a distribution $Q(x)$ over experts relative to a uniform prior $P$, the PAC-Bayesian generalization gap is bounded by the Kullback-Leibler (KL) divergence $\mathbb{E}_{x}[D_{\text{KL}}(Q(x) \| P)]$. If the prior $P$ is designed asymmetrically to place lower probability on experts with massive task-vector norms (representing high prior uncertainty or risk of parameter distortion), minimizing this KL divergence naturally yields a geometry-aware, asymmetric penalty that mathematically mirrors the SR3 objective, showing that our geometry-aware scaling is a universal principle across both deterministic and stochastic learning-theory paradigms.
    \item \textbf{Lipschitz Constant Estimation and Hyperparameter Absorption:} In practice, computing the exact global network Lipschitz constant $L_{\text{net}}$ for a large transformer model is NP-hard. However, our formulation does not require explicit knowledge of $L_{\text{net}}$. The Lagrange multiplier $\lambda$ in our objective function absorbs $L_{\text{net}}$ entirely, along with other global constant multipliers (such as $\sqrt{2K/n}$). Consequently, practitioners only need to compute the relative ratios of the task-vector norms $\|V_k\|$ across experts, and tune the single scalar hyperparameter $\lambda$ using standard cross-validation or calibration data sweeps.
    \item \textbf{Gating Entropy Regularization as a Complementary Constraint:} In Mixture-of-Experts (MoE) literature, it is standard to incorporate an entropy-based regularizer $\mathcal{L}_{\text{entropy}}(\alpha) = \beta \sum_{b=1}^B \sum_{k=1}^K \alpha_{k}(b) \log \alpha_{k}(b)$ to balance expert load or sharpen task selection. Under our Rademacher generalization framework, gating entropy represents a constraint on the effective dimension of the routing prediction simplex. While SR3 bounds the router's parameter-space capacity $\|W_k\|_2$, an entropy penalty limits the representation-space entropy of the predicted routing coefficients $\alpha_k(x)$. Specifically, high gating entropy spreads routing across multiple experts, which increases parameter-space ensembling and representation smoothing, while low gating entropy forces sparse, single-expert selection. This shows that gating entropy and SR3 are highly complementary: SR3 controls the parameter-space sensitivity of each individual routing direction, while gating entropy controls the input-dependent selection profile over the routing simplex, opening a promising direction for joint statistical-and-entropy-constrained dynamic model merging.
\end{enumerate}

Theorem~\ref{thm:rademacher} provides a rigorous mathematical connection between parameter-space task complexity ($\|V_k\|_F$) and the router's generalization capacity. It demonstrates that the contribution of each expert to the overall generalization gap is scaled by $\|V_k\|_F$. Thus, if an expert $k$ lies far from the base model in parameter space, even a small variation in its routing parameters $W_k, B_k$ can cause massive representation shifts, inflating the generalization error.

\subsection{Derivation of the Geometry-Aware Regularizer}
Having bounded the generalization complexity of our merged hypothesis class, we seek to design a regularization scheme that directly controls this complexity. 

While directly minimizing the linear Rademacher bound $\sum_{k=1}^K v_k w_k$ (where $v_k = \|V_k\|_F$ and $w_k = \sqrt{\|W_k\|_2^2 + B_k^2}$) corresponds to a non-differentiable $L_1$-like group lasso penalty at the origin, we instead define our parameter capacity constraint quadratically via a smooth, convex, ellipsoidal capacity constraint. This serves as a differentiable surrogate that preserves the asymmetric, geometry-aware scaling of each expert's parameter contribution.

\begin{theorem}[SR3 Regularization under Weighted Parameter Capacity Constraint]
Let $w_k = \sqrt{\|W_k\|_2^2 + B_k^2}$ be the norm of the routing parameters for the $k$-th expert, and let $v_k = \|V_k\|_F$ be the Frobenius norm of its corresponding task vector. To minimize the empirical calibration loss $\mathcal{L}_{\text{CE}}(w)$ subject to a smooth, geometry-aware capacity constraint on the weighted sum of routing parameter capacities:
\begin{equation}
    \sum_{k=1}^K v_k w_k^2 \le C_0'
\end{equation}
for some capacity budget $C_0' > 0$, the optimal unconstrained Lagrange objective is obtained by applying the regularizer:
\begin{equation}
    \mathcal{L}_{\text{reg}} = \lambda \sum_{k=1}^K v_k w_k^2
\end{equation}
where $\lambda \ge 0$ is a Lagrange multiplier.
\end{theorem}

\begin{proof}
Let $\mathcal{L}_{\text{CE}}(w)$ denote the empirical calibration loss over the sparse calibration split. Under statistical learning theory, to prevent generalization collapse, we must minimize the empirical risk subject to a constraint on the generalization gap, which we bound via the empirical Rademacher complexity $\mathcal{R}_n(\mathcal{H}_{\text{merged}})$.

To construct a smooth, differentiable, and convex regularization boundary at the origin, we bound the weighted capacity contribution of each expert, where each expert's routing parameter capacity $w_k^2$ is weighted by its parameter-space task-vector norm $v_k = \|V_k\|_F$:
\begin{equation}
    \sum_{k=1}^K v_k w_k^2 \le C_0'
\end{equation}
This inequality constraint models the geometry-aware capacity limit of the dynamic router. We formulate the corresponding constrained optimization problem:
\begin{equation}
    \min_{w} \mathcal{L}_{\text{CE}}(w) \quad \text{subject to} \quad \sum_{k=1}^K v_k w_k^2 - C_0' \le 0
\end{equation}
Using the method of Lagrange multipliers, we construct the Lagrangian function with multiplier $\lambda \ge 0$:
\begin{equation}
    \mathcal{L}_{\text{Lag}}(w, \lambda) = \mathcal{L}_{\text{CE}}(w) + \lambda \left( \sum_{k=1}^K v_k w_k^2 - C_0' \right)
\end{equation}
Since $C_0'$ is a constant, optimizing the Lagrangian over the routing parameters $w$ is equivalent to minimizing:
\begin{equation}
    \min_{w} \mathcal{L}_{\text{CE}}(w) + \lambda \sum_{k=1}^K v_k w_k^2
\end{equation}
which is precisely our empirical objective with the regularization term:
\begin{equation}
    \mathcal{L}_{\text{reg}} = \lambda \sum_{k=1}^K v_k w_k^2
\end{equation}
This completed regularizer scales the penalty on the squared routing parameter norm $w_k^2$ by the linear task-vector norm $v_k$.

From an optimization perspective, this asymmetric scaling ensures that during gradient descent, the effective weight decay force acting on the routing parameters of expert $k$ is directly proportional to its task-vector norm $v_k$. Since experts with large task-vector norms $v_k$ represent highly sensitive directions in parameter space (where even a small change in routing coefficients causes massive representation distortion), SR3 penalizes their parameters more aggressively, forcing their routing weights to remain small and conservative. Conversely, low-norm experts (which have low generalization sensitivity) are lightly penalized, preserving high routing flexibility on simpler domains. This linear scaling of the regularizer with respect to the task vector norm ($v_k$) is theoretically consistent with Theorem~\ref{thm:rademacher} and prevents the extreme over-regularization of high-complexity experts that occurs under a squared scaling factor ($v_k^2$), achieving an optimal balance across all task domains during calibration.
\end{proof}

This provides an airtight, first-principles proof that **scaling weight decay by the task-vector norm is the theoretically optimal strategy to bound generalization error in weight-space model merging**.

\subsection{The SR3 Loss Objective}
The finalized \textbf{Spectral and Rademacher-guided Routing Regularization (SR3)} penalty is formulated layer-wise as:
\begin{equation}
    \mathcal{L}_{\text{SR3}} = \lambda_{\text{SR3}} \sum_{l=1}^L \sum_{k=1}^K \Gamma_k^{(l)} \left( \|W_{l, k}\|_2^2 + B_{l, k}^2 \right)
\end{equation}
where the expert-specific scaling multiplier $\Gamma_k^{(l)}$ is defined according to the parameter-space norm variant:
\begin{itemize}
    \item \textbf{Frobenius Variant (SR3-F):} Scales the penalty by the linear Frobenius norm of the task vector:
    \begin{equation}
        \Gamma_k^{(l)} = \|V_k^{(l)}\|_F
    \end{equation}
    \item \textbf{Spectral/Operator Variant (SR3-S):} Scales the penalty by the linear operator (spectral) norm (i.e., the largest singular value) of the task vector:
    \begin{equation}
        \Gamma_k^{(l)} = \|V_k^{(l)}\|_{op} = \sigma_{\max}\left(V_k^{(l)}\right)
    \end{equation}
\end{itemize}

While the Frobenius norm bounds the total accumulated change across all parameter dimensions, the operator norm specifically captures the \emph{worst-case representation distortion} under the linear transformation $V_k^{(l)}$. Under a tighter lipschitz analysis where we consider layer-wise representation bounds, the operator norm provides a tighter generalization guarantee. Empirically, we show that the spectral variant SR3-S indeed yields superior calibration performance, confirming that constraining worst-case transformation distortion is a tighter generalization constraint than constraining the Frobenius norm.

\paragraph{Computational Complexity and Scalability.} A potential concern for the spectral norm variant (SR3-S) is the cost of Singular Value Decomposition (SVD) for high-dimensional matrices, which scales as $\mathcal{O}(D^3)$ for an expert layer of dimension $D$. However, because the task vectors $V_k^{(l)}$ are static and pre-computed once offline prior to router training, this profiling step introduces zero runtime overhead during calibration. Furthermore, to scale this offline step to giant models, we do not need to perform exact SVD; instead, we can compute the spectral norm efficiently using power iterations or randomized SVD. Power iterations require only simple matrix-vector multiplications, reducing the computational complexity to $\mathcal{O}(m \cdot D^2)$ for $m$ iterations (typically $m \le 3$ is sufficient). This reduces the complexity by a factor of $D$ (over three orders of magnitude for $D=4096$), making the profiling step extremely fast and highly scalable (see Section~\ref{sec:conclusion} and Appendix~\ref{app:power_iteration} for empirical validation).

The complete multi-objective calibration loss optimized during the router training phase is:
\begin{equation}
    \mathcal{L}_{\text{total}} = \mathcal{L}_{\text{CE}} + \mathcal{L}_{\text{SR3}}
\end{equation}
where $\mathcal{L}_{\text{CE}}$ is the task classification Cross-Entropy loss computed over the calibration split.

\subsection{Direct Rademacher Minimization: Smoothed $L_1$ Group-Lasso Variant}
To directly minimize the linear Rademacher complexity bound derived in Theorem~\ref{thm:rademacher} (rather than using the quadratic ellipsoidal surrogate), we also propose and evaluate a smoothed $L_1$ Group-Lasso variant of our regularizer. While a pure $L_1$ group-lasso penalty of the form $\sum_k v_k w_k$ is non-differentiable at the origin, we introduce a small smoothing parameter $\epsilon_{\text{smooth}} > 0$ to guarantee differentiability and numerical stability during standard backpropagation:
\begin{equation}
    \mathcal{L}_{\text{SR3-L1}} = \lambda_{\text{SR3-L1}} \sum_{l=1}^L \sum_{k=1}^K \Gamma_k^{(l)} \sqrt{\|W_{l, k}\|_2^2 + B_{l, k}^2 + \epsilon_{\text{smooth}}}
\end{equation}
where $\Gamma_k^{(l)}$ is again either the linear Frobenius norm $\|V_k^{(l)}\|_F$ (giving the \textbf{SR3-F-L1} variant) or the spectral norm $\|V_k^{(l)}\|_{op}$ (giving the \textbf{SR3-S-L1} variant). In our implementation, we set $\epsilon_{\text{smooth}} = 10^{-8}$. This formulation represents the exact differentiable minimizer of the derived Rademacher bound, avoiding any surrogate approximations and providing a direct, unconstrained empirical realization of statistical learning theory.

\subsection{Regularization Scheduling: Resolving the $L_1$ Group-Lasso Paradox}
Although the smoothed $L_1$ Group-Lasso penalty $\mathcal{L}_{\text{SR3-L1}}$ directly implements the linear Rademacher bound, it faces an optimization barrier near the origin (described in Section~\ref{sec:l1_paradox}), where the extremely steep early gradients can over-repress the routing weights of complex tasks and prevent them from learning. To bypass this optimization barrier, we propose a dynamic \textbf{Regularization Scheduling} (or warm-up) scheme. We define a time-varying regularizer that starts as the smooth, vanishing-gradient quadratic surrogate $\mathcal{L}_{\text{SR3}}$ at the beginning of calibration and smoothly transitions to the theoretically optimal direct $L_1$ penalty $\mathcal{L}_{\text{SR3-L1}}$ as training progresses. Specifically, at training epoch $t \in \{0, \dots, T-1\}$, we optimize:
\begin{equation}
    \mathcal{L}_{\text{SR3-Sched}}(t) = \left(1 - \frac{t}{T-1}\right) \mathcal{L}_{\text{SR3}} + \left(\frac{t}{T-1}\right) \mathcal{L}_{\text{SR3-L1}}
\end{equation}
This schedule is applied to both the Frobenius norm variant (\textbf{SR3-F-L1-Sched}) and the Spectral norm variant (\textbf{SR3-S-L1-Sched}). By utilizing the smooth quadratic penalty early in training, the routing weights are allowed to escape from zero and fit the calibration classification loss. As the weights grow, the regularizer smoothly morphs into the direct $L_1$ penalty, enforcing the tighter, learning-theoretic generalization constraint by the end of training. This achieves a perfect marriage of optimization friendliness and statistical learning optimality.

